% ==========================================================================
% GENERATED by scripts/split_paper.py. Do not edit this file.
%
% Source: imported_papers/grokking/iclr2027_conference_v2.tex
% Part 13: Scope and limits of the measures
%
% Edits belong upstream, in the published paper. To change what the
% chapter includes, change imports/grokking.toml and re-run
% the splitter.
% ==========================================================================
\section{Scope and limits of the measures}
The paper leans on a small set of measures, and each does a different job. This section says where each stops being informative.

\subsection{Participation ratio versus spectral entropy}
\label{app:pr}
Spectral entropy depends on the basis it is computed in, and on multiplication that dependence breaks it completely: every rule reads near $1.0$, including backpropagation, whose circuit is the most concentrated by every unit level measure (Table~\ref{tab:prent}). The reason is a fact about characters. A solution of $a \cdot b$ is built from the multiplicative characters of $\mathbb{Z}_p^\times$, and written in our additive basis every one of them has the same magnitude $\sqrt{p}$ at every frequency \citep{ireland1990classical}, so even a circuit concentrated on a single character reads as perfectly flat.

Participation ratio does not depend on a choice of basis, separates the rules on every operation, and does not move when the readout is rescaled, since every outgoing norm scales by the same factor. It therefore leads throughout, and spectral entropy appears only on addition, where the additive basis is the right one and the two measures agree.

\begin{table}[!ht]
\caption{Participation ratio separates the rules on every operation, while spectral entropy saturates on multiplication. Spectral entropy in the additive Fourier basis versus participation ratio, by operation ($p=31$, three seeds). Entropy separates the rules on addition and subtraction but saturates for every rule on multiplication, as the Gauss sum argument predicts. Participation ratio separates the rules everywhere.}
\label{tab:prent}
\begin{center}
\small
\begin{tabular}{lcccccc}
\toprule
 & \multicolumn{3}{c}{Spectral entropy} & \multicolumn{3}{c}{Participation ratio} \\
Operation & backprop & knob & DFC & backprop & knob & DFC \\
\midrule
Addition       & $0.85$ & $0.97$ & $1.00$ & $0.60$ & $0.99$ & $0.92$ \\
Subtraction    & $0.85$ & $0.96$ & $1.00$ & $0.59$ & $0.99$ & $0.93$ \\
Multiplication & $1.00$ & $1.00$ & $1.00$ & $0.43$ & $0.99$ & $0.88$ \\
\bottomrule
\end{tabular}
\end{center}
\end{table}

\subsection{Checks that the algorithm is shared}
\label{app:algo}
The algorithm in question is the Fourier computation the grokking literature derives for modular addition. The embedding maps each operand onto a few frequencies $k$. The hidden layer combines them into $\cos$ and $\sin$ of $k(a+b)$. The logit for a candidate answer $c$ is then a sum of $\cos(2\pi k(a+b-c)/p)$ over those frequencies, which peaks when $c$ equals $a+b$.

Existing theory derives this circuit from the task and the architecture, with no reference to how credit is assigned \citep{nanda2023progress,chughtai2023toy,gromov2023grokking,morwani2024feature}. The following checks confirm that both rules compute it and that nothing left over is a second algorithm.

\paragraph{Frequency ablation.} Keeping only the top input frequencies recovers $95\%$ of full accuracy from $6.7$ frequencies for backpropagation and $9.3$ for DFC, so both rest on a periodic Fourier code.

\paragraph{Fit to the trigonometric basis.} Regressing the full logit tensor onto the $(a+b-c)$ trigonometric basis, the form a modular addition circuit must take, explains $0.90$ of DFC's logit variance against $0.57$ of backpropagation's. DFC realizes the circuit cleanly and backpropagation more noisily, consistent with backpropagation's lower per-neuron periodicity. The unexplained share is noise rather than a second algorithm: backpropagation's residual agrees with itself across initializations at only $0.235$, so it does not reproduce even within one rule, where an algorithm would.

\paragraph{Agreement across runs, at the logits and in the hidden layer.} Two runs of one rule share an algorithm by construction, so their agreement is the ceiling any pair of rules can reach, and with the train and test split fixed a seed changes only the initialization. Backpropagation agrees with itself on the full logit tensor at correlation $0.446$, against a shuffled floor of $0.000$, and agreement between different rules falls in the same range, $0.427$ to $0.467$. The same holds one level below the logits, where linear CKA between first hidden layer activations averages $0.69$ between backpropagation and the knob at $\beta=1$, against seed to seed agreement of $0.57$ within backpropagation and $0.74$ within the knob (three seeds each, shared data split). Changing the rule costs no more agreement, functional or representational, than changing the initialization.

\paragraph{The one published alternative.} Regressing the same logit tensors onto the pizza basis of \citet{zhong2023clock}, $\lvert\cos(\pi k(a-b)/p)\rvert \cdot \cos(\pi k(a+b-2c)/p)$, explains $0.002$ to $0.005$ of the variance for every rule, while the same estimator scores $1.000$ on a synthetic pizza tensor. No rule here implements it, so no second algorithm hides behind the trigonometric fit.

\subsection{Output scale and weight norm}
\label{app:norm}
Raw parameter norm is only meaningful relative to the scale of the network's outputs, since rescaling the readout rescales the margin with it. The two rules settle at different scales, DFC's margin near $0.54$ against backpropagation's $5.0$, so their raw norms, $13$ against $26$, cannot be compared. The scale invariant quantity is the margin normalized norm $\|\theta\| / \gamma$, the notion of efficiency the grokking literature uses \citep{varma2023circuit}, and under it backpropagation is the more efficient solution by a factor of about $4.6$ (Table~\ref{tab:norms}). Two consequences follow. First, the distributed circuit is not a minimum norm solution, and we claim none. Second, robustness and norm efficiency are separate axes: the knob is as robust as DFC under dropout while sharing backpropagation's norm, margin and efficiency.

\begin{table}[!ht]
\caption{Raw norm favors DFC, the scale invariant ratio favors backpropagation, and the credit knob shows the two are independent. Weight norm, training margin, and margin normalized norm at matched weight decay (modular addition, $p=31$, mean and standard deviation over four seeds, all at test accuracy $\geq 0.97$).}
\label{tab:norms}
\begin{center}
\begin{tabular}{lccc}
\toprule
Rule & $\|\theta\|$ & margin $\gamma$ & $\|\theta\| / \gamma$ \\
\midrule
Backpropagation    & $26.10 \pm 0.15$ & $5.02 \pm 0.02$ & $5.20 \pm 0.04$ \\
Knob ($\beta{=}1$) & $26.65 \pm 0.10$ & $5.15 \pm 0.03$ & $5.17 \pm 0.02$ \\
DFC                & $13.09 \pm 0.02$ & $0.54 \pm 0.00$ & $24.20 \pm 0.06$ \\
\bottomrule
\end{tabular}
\end{center}
\end{table}

\subsection{Spread is not neuron purity}
\label{app:dissoc}
Max-margin theory describes the grokked solution by two structural properties at once, neurons that each tune to a single frequency and a circuit spread across many units \citep{morwani2024feature}, and the two dissociate (Table~\ref{tab:plane}). Backpropagation reaches neither, at participation $0.595$ and a single frequency fraction of $0.349$. The knob reaches the spread and not the purity, at $0.992$ and $0.487$. Nero reaches the spread, $0.960$, while its purity falls to $0.255$, below backpropagation's own, and only DFC approaches both, at $0.923$ and $0.560$. The same table carries the minimal breaking set of Section~\ref{sec:fragile}, defined in the caption, and it rises with participation across every condition.

\begin{table}[!ht]
\caption{The max-margin plane and the minimal breaking set (modular addition, $p=31$,
three seeds). These are an independent run from Table~\ref{tab:twoaxis}, so the rules they share differ in their per-run values. Single frequency fraction is the share of units whose spectrum is dominated by one frequency, and $\#$freqs the number of distinct dominant frequencies. The break set is the smallest set a greedy accuracy search must delete to drop test below $0.5$, and norm pct the mean norm percentile of the units it chooses. Nero and LARS are the validated implementations of Appendices~\ref{app:named} and~\ref{app:zoo}.}
\label{tab:plane}
\begin{center}
\small
\begin{tabular}{lcccccc}
\toprule
Rule & Test & Circuit PR & Single-freq & $\#$freqs & Break set & Norm pct \\
\midrule
Backpropagation    & $0.998$ & $0.595\,(0.023)$ & $0.349\,(0.033)$ & $15.0$ & $23.3$ & $0.14$ \\
LARS               & $0.989$ & $0.530\,(0.029)$ & $0.309\,(0.055)$ & $13.7$ & $16.3$ & $0.09$ \\
Knob ($\beta{=}0.25$) & $0.997$ & $0.940\,(0.017)$ & $0.439\,(0.056)$ & $12.3$ & $43.0$ & $0.26$ \\
Knob ($\beta{=}0.5$)  & $0.986$ & $0.975\,(0.002)$ & $0.469\,(0.039)$ & $13.3$ & $59.0$ & $0.34$ \\
Knob ($\beta{=}1$)    & $0.976$ & $0.992\,(0.002)$ & $0.487\,(0.016)$ & $13.3$ & $60.7$ & $0.39$ \\
Nero               & $0.824$ & $0.960\,(0.003)$ & $0.255\,(0.046)$ & $15.0$ & $37.0$ & $0.30$ \\
DFC                & $1.000$ & $0.923\,(0.003)$ & $0.560\,(0.019)$ & $15.0$ & $53.3$ & $0.31$ \\
\bottomrule
\end{tabular}
\end{center}
\end{table}
